\chapter*{Presentación}
\markboth{Presentación}{Presentación}
\addcontentsline{toc}{chapter}{Presentación}

Este temario organiza un primer curso riguroso de álgebra lineal alrededor de los espacios vectoriales y las transformaciones lineales. La ruta empieza con campos y subespacios, pasa por independencia, bases y dimensión, y llega a la representación matricial y los espacios cociente.

La columna vertebral es
\[
  \text{espacios}\longrightarrow\text{bases}\longrightarrow
  \text{transformaciones}\longrightarrow\text{rango--nulidad}\longrightarrow
  \text{sistemas}\longrightarrow\text{matrices}\longrightarrow\text{cocientes}.
\]
Cada capítulo delimita los conceptos y resultados que conviene estudiar y demostrar. Se presupone familiaridad con operaciones elementales y sistemas lineales; cuando sea útil, se revisarán esos procedimientos desde su fundamento vectorial.

Para trabajar el curso en dos niveles, se incluyen las dos ediciones de Serge Lang: \emph{Introduction to Linear Algebra}, segunda edición, y \emph{Linear Algebra}, tercera edición. La primera sirve como texto conciso de entrada; la tercera reorganiza y amplía varios temas. Hoffman y Kunze ofrecen un desarrollo axiomático clásico, mientras que Axler aporta una perspectiva centrada en operadores y espacios vectoriales \cite{lang1986,lang1987,hoffmankunze1971,axler2024}.
